\documentclass[12pt,a4paper]{article}

\usepackage[T2A]{fontenc}
\usepackage[utf8]{inputenc}
\usepackage[russian]{babel}
\usepackage{geometry}
\usepackage{setspace}
\usepackage{enumitem}
\usepackage{hyperref}
\usepackage{verbatim}
\usepackage{parskip}

\geometry{left=25mm,right=25mm,top=20mm,bottom=20mm}
\onehalfspacing
\setlist[itemize]{leftmargin=1.25cm}
\setlist[enumerate]{leftmargin=1.25cm}

\begin{document}

\begin{center}
    \textbf{\Large Архитектура продукта}\\[0.5em]
    \textbf{\large Компьютерная игра ``The Hero''}
\end{center}


\section*{1. Архитектурно значимые особенности продукта}

Архитектурно значимое ядро продукта задаётся игровым циклом <<карта $\rightarrow$ событие/бой $\rightarrow$ награда $\rightarrow$ усиление армии/базы>>. Пользователь перемещает героя по клеточной карте, взаимодействует с объектами, вступает в бой с нейтральными противниками, получает ресурсы, опыт и дополнительные награды, после чего усиливает армию через базу и постройки. Поэтому архитектура должна строиться не вокруг отдельных экранов, а вокруг единого состояния игры, последовательно проходящего через режим карты, боя, базы/найма, информационных окон и настроек.

Карта мира представляет собой клеточную тайловую сетку, по которой герой перемещается автоматически по рассчитанному маршруту после выбора цели. На карте размещаются ресурсы, нейтральные противники, статичные противники-отряды, база и нейтральные строения. Это требует наличия модели карты, модели объектов карты, механизма маршрутизации, модуля обработки достижения цели и механизма фиксации состояния объектов.

Бой выполняется как отдельный режим, не связанный с пространственной моделью карты. Боевая подсистема получает состав сторон, параметры юнитов, инициативу и выбранные цели, выполняет расчёт столкновения и возвращает результат в виде потерь, исхода боя и награды. Такая схема требует самостоятельного боевого модуля с собственной моделью состояния, отделённой от карты и маршрутизации.

База и найм образуют отдельный архитектурный контур. Внутри него должны храниться постройки, уровни построек, доступные типы юнитов, накопленный объём найма, правила пополнения по недельному циклу и операции улучшения. Это требует выделения модуля базы, модуля найма, модуля недельного цикла и механизма улучшения уже нанятых юнитов.

Подсистема сохранения должна обрабатывать как основное игровое состояние, так и связанные с ним метаданные: день и неделю, состояние объектов карты, состав армии, состояние базы и пользовательские настройки. Для этого она должна быть самостоятельной, версионируемой и устойчивой к частичным повреждениям данных.

Пользовательский интерфейс также является полноценной частью архитектуры. Он включает главное меню, экран карты, экран боя, экран базы/найма, окно настроек, информационные окна, локализацию интерфейса, управление звуком, выбор сложности и раздел помощи или обучения. Следовательно, слой представления должен иметь собственную структуру экранов и окон, но не содержать игровую логику.

\section*{2. Архитектурный подход}

Для данного продукта обоснован модульный монолит внутри Unity. Все режимы работают внутри одного клиентского приложения, используют общее игровое состояние и не требуют разделения на сетевые сервисы, серверную часть или внешние хранилища.

Базовый принцип архитектуры состоит в том, что единое доменное состояние игры находится в центре, а режимы карты, боя, базы, настроек и сохранения работают как согласованные подсистемы над этим состоянием. Это позволяет централизованно изменять состояние героя, армии, ресурсов, карты, базы и временного цикла без дублирования логики между экранами.

Внутри такого монолита целесообразно выделить четыре слоя. \textbf{Первый слой} представляет собой доменную модель игры и включает основные сущности предметной области: \textbf{героя} (идентификатор, тип или имя, характеристики, текущее состояние, положение на карте, связанная армия), \textbf{армию} (состав, типы и количество юнитов, владелец, текущее состояние), \textbf{юнита} (тип, боевые параметры, текущее состояние, принадлежность стороне), \textbf{ресурсы} (виды ресурсов, текущее количество, принадлежность, ограничения на расход), \textbf{карту} (области или точки, связи между ними, размещённые объекты, текущее состояние элементов карты), \textbf{базу} (идентификатор, положение на карте, доступные функции, связанные ресурсы, текущее состояние), \textbf{награду} (тип, величина, источник, условие получения, статус выдачи). \textbf{Второй слой} образуют игровые подсистемы: \textbf{подсистема карты}, отвечающая за перемещение героя и взаимодействие с объектами; \textbf{подсистема боя}, выполняющая инициацию столкновения, определение участников, расчёт результата и применение последствий; \textbf{подсистема экономики}, обеспечивающая начисление, списание и проверку достаточности ресурсов; \textbf{подсистема найма}, проверяющая доступность найма, стоимость и добавление юнитов в армию; \textbf{подсистема наград и результатов событий}, применяющая последствия боёв и взаимодействий с объектами; \textbf{подсистема сохранения}, собирающая и восстанавливающая игровое состояние. \textbf{Третий слой} составляет слой представления, включающий \textbf{игровые экраны}, \textbf{окна} и прочие элементы \textbf{пользовательского интерфейса}. \textbf{Четвёртый слой} образует инфраструктурный уровень, в который входят \textbf{сериализация JSON}, \textbf{загрузка конфигурации}, \textbf{локализация}, \textbf{аудиоподсистема} и \textbf{файловый ввод-вывод}. Такое разделение позволяет изолировать данные предметной области, игровую логику, интерфейс и технические механизмы поддержки приложения.

Боевая подсистема должна быть отделена от подсистемы карты. Она получает на вход состав сторон, параметры юнитов и иные числовые характеристики, необходимые для расчёта столкновения, а на выходе возвращает результат боя, потери и награды. Пространственная логика боя, включая поле, клетки и позиционирование, в архитектуру боевого модуля не включается.

\section*{3. Компоненты и модули}

\subsection*{3.1. Ядро приложения}

\textbf{Модуль управления приложением и сессией.}

Отвечает за запуск игры, переходы из главного меню, создание новой игры, продолжение сохранённой сессии и завершение работы приложения. Этот модуль координирует вход в игру и переключение между основными режимами.

\textbf{Модуль состояния игры.}

Хранит текущее состояние игры как единый агрегат: состояние героя, армии, ресурсов, карты, базы, игровых дней и недель, выбранной сложности, пользовательских настроек, флагов посещения объектов и результатов боёв. Это центральный модуль, потому что все игровые подсистемы работают с одной и той же моделью данных.

\subsection*{3.2. Игровые подсистемы}

\textbf{Модуль карты мира.}

Отвечает за модель тайловой карты, координаты, доступность клеток, размещение объектов карты, отображение точки назначения, расчёт маршрута, расход очков движения и завершение перемещения. Он формирует основной режим перемещения по миру и подготавливает переход к сценариям взаимодействия.

\textbf{Модуль объектов карты и взаимодействий.}

Обрабатывает достижение выбранного объекта и запускает нужный сценарий: сбор ресурса, переход в бой, выдачу награды, открытие базы, захват или пометку нейтрального строения как использованного либо захваченного. Этот модуль объединяет различные типы объектов карты в единый механизм взаимодействия.

\textbf{Модуль героя и армии.}

Хранит параметры героя, состав армии, слоты армии, типы и количество юнитов в каждом отряде, а также ограничения на размер армии. Он выделяется отдельно, потому что армия участвует в карте, базе, бою и сохранении.

\textbf{Модуль юнитов и характеристик.}

Содержит определения типов юнитов и их параметров: здоровье или численность, атаку, защиту или броню, урон, инициативу или скорость. Он отделяется от текущего состояния армии, потому что один и тот же тип юнита может использоваться в разных армиях, в найме, в улучшениях и в расчёте боя.

\textbf{Боевой модуль.}

Включает формирование сторон, очередь ходов по инициативе, приём выбора целей игроком, расчёт урона, учёт потерь, определение победы или поражения и формирование результата боя. Это самостоятельный модуль, поскольку бой выполняется в отдельном режиме и должен быть воспроизводим при одинаковых исходных данных.

\textbf{Модуль экономики и ресурсов.}

Учитывает получение ресурсов на карте и после боёв, расход на найм, улучшения построек и улучшение ранее нанятых юнитов. Он нужен как отдельный модуль, потому что ресурсные операции должны проверяться и применяться единообразно во всех сценариях.

\textbf{Модуль базы, построек и найма.}

Хранит перечень построек базы, их уровень, доступные типы юнитов, накопленный доступный найм, правила пополнения раз в 7 ходов и операции найма или улучшения. Он выделяется отдельно, поскольку база имеет собственный жизненный цикл и собственные данные, не совпадающие ни с картой, ни с боем.

\textbf{Модуль наград и результатов событий.}

Формирует и применяет последствия побед и событий: ресурсы, опыт героя, артефакты и другие изменения параметров, если они поддерживаются механикой. Этот модуль связывает результат взаимодействия с последующим изменением общего состояния игры.

\textbf{Модуль хода, игрового дня и недели.}

Инкрементирует глобальный ход, завершает игровой день по кнопке <<Завершить ход>>, инициирует автосохранение и пополнение найма по игровой неделе. Этот модуль координирует временной цикл и связанные с ним системные изменения.

\subsection*{3.3. Подсистемы представления и инфраструктуры}

\textbf{UI-модуль экранов и окон.}

Включает главное меню, экран карты, экран боя, экран базы и найма, окно настроек, окна информации о герое, армии, ресурсах, боях, событиях и постройках. Он отвечает за отображение состояния системы и передачу пользовательских команд в игровые подсистемы.

\textbf{Модуль пользовательских настроек.}

Отвечает за громкость, включение и выключение звука, выбор языка интерфейса, уровень сложности, а также настройку скорости анимаций или их отключение в поддерживаемых режимах. Он изолирует пользовательские параметры от игрового состояния.

\textbf{Модуль локализации.}

Обеспечивает русский и английский интерфейс для UI-элементов. Он должен быть самостоятельным, потому что текстовое представление интерфейса используется на всех основных экранах.

\textbf{Модуль помощи и обучения.}

Реализует раздел помощи и обучающий режим либо набор подсказок для первых уровней и первых боёв. Он отвечает за подачу вспомогательной информации, не вмешиваясь в базовую игровую логику.

\textbf{Модуль сохранения и загрузки.}

Отвечает за сериализацию состояния игры, запись ручных и автоматических сохранений, хранение резервной копии, валидацию структуры при чтении, проверку версии формата и обработку ошибок с безопасным восстановлением. Он является отдельной подсистемой из-за собственной логики целостности данных.

\textbf{Модуль конфигурационных данных.}

Загружает параметры юнитов, построек, наград, составов противников и другие конфигурации баланса из JSON. Он отделяет изменяемые игровые параметры от исполняемой логики.

\textbf{Аудиомодуль.}

Обеспечивает проигрывание музыки и звуковых эффектов, а также применение пользовательских настроек громкости и отключения звука. Он инкапсулирует всю работу со звуковой подсистемой приложения.

\section*{4. Связи между модулями}

Управление приложением и UI верхнего уровня инициируют создание новой игры или загрузку сохранения. После этого создаётся или восстанавливается единое состояние игры, которое передаётся в модуль карты как основной режим прохождения.

Модуль карты использует модуль объектов карты и взаимодействий. Когда герой достигает цели, обработчик объекта определяет тип результата: мгновенный сбор ресурса, открытие базы, переход в бой, выдача награды, захват нейтрального строения. Сам модуль карты не содержит логику боя или найма; он только инициирует соответствующий сценарий.

При запуске боя модуль взаимодействий передаёт в боевой модуль текущую армию игрока из модуля героя и армии, а также выбранный вражеский отряд с его параметрами из конфигурации и состояния карты. Боевой модуль использует определения юнитов и возвращает результат: потери сторон, исход боя, награду и обновлённое состояние армии. Затем модуль наград и результатов событий применяет последствия к общему состоянию игры.

Модуль базы, построек и найма получает доступ к модулю экономики и ресурсов для проверки стоимости, к модулю героя и армии для добавления юнитов в слоты и к модулю хода и недели для расчёта накопленного найма. Таким образом, база не существует изолированно: она зависит от времени, ресурсов и текущего состава армии.

Модуль хода и недели воздействует сразу на несколько подсистем: обновляет счётчики времени, инициирует недельное пополнение найма в базе, запускает автосохранение и фиксирует завершение игрового дня. Это один из ключевых координационных модулей всей архитектуры.

Модуль сохранения и загрузки работает не с отдельными экранами, а с единым состоянием игры и пользовательскими настройками. При сохранении он получает текущее состояние игры. При загрузке он выполняет валидацию, проверку версии и либо восстанавливает корректное состояние, либо предлагает резервное сохранение или запуск новой игры.

UI-модуль связан со всеми игровыми подсистемами только через чтение их состояния и вызов разрешённых действий. Интерфейс отображает данные и передаёт команды, но не хранит собственную игровую логику, что обеспечивает консистентное сохранение и корректные переходы между режимами.

\section*{5. Данные и сущности}

Ключевой агрегат данных --- \textbf{состояние игры}. В него входят состояние героя, состав армии, ресурсы, состояние карты, состояние базы, текущий игровой день и неделя, результаты посещения объектов, выбранные параметры сложности и данные, необходимые для сохранения и восстановления игровой сессии.

\textbf{Карта} представляется как набор клеток и размещённых на них объектов. Для объектов необходимы по меньшей мере тип, координаты и состояние. Состояние объекта должно позволять различать собранный ресурс, побеждённого противника, посещённый объект и захваченное строение.

\textbf{Герой} является сущностью с параметрами, участвующей в перемещении по карте и получении наград. В архитектуре он должен поддерживать изменяемые характеристики, накопление опыта и связь с армией.

\textbf{Армия} хранится как набор слотов, где каждый слот содержит один отряд. \textbf{Отряд} представляет собой тип юнита и его текущее количество. Такое представление позволяет фиксировать состав армии, потери после боя и изменения после найма или улучшения.

\textbf{Тип юнита} хранит боевые характеристики: здоровье или численность, атаку, защиту или броню, урон, инициативу или скорость. Отдельно от него хранится \textbf{экземпляр отряда}, содержащий текущее количество, чтобы после боя можно было зафиксировать выживших без автоматического восстановления.

\textbf{База} включает перечень построек. \textbf{Постройка} содержит тип, уровень, доступный тип юнита для найма, объём доступного найма на текущий момент и данные для улучшения. Для накопительного найма у постройки должен существовать сохраняемый счётчик накопленного предложения, а не только текущая доступность на экране.

\textbf{Ресурсы} хранятся в виде отдельной сущности-кошелька, поскольку используются в нескольких независимых сценариях: сбор на карте, награда после боя, найм, улучшение построек и улучшение ранее нанятых юнитов.

\textbf{Награда} выделяется в отдельную сущность, объединяющую ресурсы, опыт героя, артефакты и иные результаты события или боя. Это позволяет единообразно применять последствия различных игровых сценариев.

\textbf{Сохранение} должно содержать данные о версии формата, обязательных полях, текущем состоянии игры и ссылочно независимую структуру, пригодную для проверки целостности и миграции в пределах одной мажорной версии.

\textbf{Пользовательские настройки} отделяются от основного игрового состояния и включают по меньшей мере звук, язык интерфейса, сложность и параметры анимации в пределах реально поддерживаемых функций.

\section*{6. Сценарии работы системы}

\subsection*{Сценарий 1. Запуск и вход в сессию}

Пользователь запускает приложение, попадает в главное меню, выбирает новую игру либо продолжение. При новой игре инициализируется начальное состояние карты, героя, армии, базы и ресурсов. При загрузке читается сохранение, проходит валидация, затем состояние восстанавливается. При ошибке загрузки пользователь получает сообщение и вариант отката к резервной копии или начала новой игры.

\subsection*{Сценарий 2. Перемещение по карте}

Пользователь выбирает клетку или объект. Модуль карты рассчитывает кратчайший доступный маршрут, списывает очки движения и перемещает героя до достижения цели или исчерпания движения. По достижении цели управление передаётся модулю взаимодействий.

\subsection*{Сценарий 3. Взаимодействие с ресурсом или объектом карты}

Если достигнут объект-ресурс, ресурс добавляется в экономику, а объект меняет состояние на собранный. Если достигнута база, открывается экран базы и найма. Если достигнут охраняемый объект или противник, запускается бой. Если достигнуто нейтральное строение, выполняется его сценарий взаимодействия, а объект помечается как использованный или захваченный в соответствии с типом действия.

\subsection*{Сценарий 4. Бой}

В боевой режим передаются армия игрока и вражеский отряд. Формируется порядок ходов по инициативе. Игрок выбирает цель для своего отряда, выполняется расчёт урона, фиксируются потери. После завершения боя определяется победа или поражение, выжившие отряды возвращаются в общее состояние игры, победителю начисляется награда. Пространственная модель боя при этом не используется.

\subsection*{Сценарий 5. Усиление армии и базы}

После посещения базы пользователь открывает список построек, видит доступный найм и уровни построек, тратит ресурсы на найм либо улучшение. При наличии улучшенной постройки возможен апгрейд ранее нанятых юнитов соответствующего типа до текущего уровня. Итоговые изменения применяются одновременно к ресурсам, состоянию базы и армии.

\subsection*{Сценарий 6. Завершение хода}

Пользователь завершает игровой ход на глобальной карте. Система фиксирует выполненные действия, обновляет состояние карты, армии и героя, увеличивает счётчик дня, проверяет смену игровой недели, при необходимости пополняет доступный найм и создаёт автосохранение с резервной копией предыдущего состояния.

\subsection*{Сценарий 7. Работа с настройками, локализацией и обучением}

Пользователь открывает настройки, меняет громкость, включает или выключает звук, выбирает язык интерфейса, сложность и, при наличии, параметры анимации. Отдельно доступны раздел помощи и обучения либо подсказки для первых боёв и уровней. Эти параметры должны применяться без нарушения основного игрового состояния.

\section*{7. Схема архитектуры}

\begin{verbatim}
[Управление приложением / Главное меню]
                |
                v
      [Единое состояние игры]
                |
    +-----------+-----------+----------------+-----------------+
    |                       |                |                 |
    v                       v                v                 v
[Карта мира]      [База / постройки / найм] [Сохранение]   [Настройки]
    |                       |                |                 |
    v                       v                v                 v
[Объекты карты]     [Экономика / ресурсы] [JSON + backup] [Звук / язык /
    |                       |                                  сложность /
    +-----------+-----------+                                  анимации]
                |
                v
        [Герой и армия]
                |
                v
        [Юниты и характеристики]
                |
                v
           [Боевой модуль]
                |
                v
      [Награды / опыт / артефакты]
                |
                v
      [Обновление общего состояния]
                |
                v
      [UI экранов и информационных окон]
\end{verbatim}

Смысл схемы состоит в том, что карта, бой, база и сохранение не существуют как независимые мини-приложения, а работают поверх одного и того же состояния игры.

\end{document}
